% -------------------------------------------------------
% ASSIGNED TO: \textcolor{green}{Person 5}
% -------------------------------------------------------
% INSTRUCTION: Paste existing Section 3.6 through 3.6.4 content here.
% Convert first-person to third person throughout. No new content required.
% -------------------------------------------------------
 
\subsubsection{Defining Quality in Symbolic Explanations: Fidelity, Comprehensibility, and Stability}
When extracting "If-Then" rules from complex black box models, the quality of the rules are usually measured 
across three conflicting dimensions:
\\\\
Fidelity - This metric measures the degree to which an explanation accurately mimics the behaviour of 
the black-box model \cite{Guidotti2021}. High fidelity translates as the extracted rules being highly faithful 
to the model’s actual decision-making and not falling into the disagreement problem (different explainability 
methods produce varying explanations for the same AI model)or just being an approximation \cite{Miro2025}.
\\\\
Comprehensibility - Is defined by the sparsity of the model or the length or number of the generated 
rules for explanation \cite{Arrieta2020,Yedjour2024}. It relates to the cognitive load required 
for humans to understand and interpret the explanation. Literature suggests that as the number of conditions 
in a rule increases, human understanding decays, a phenomenon often referred to as Rule Explosion 
\cite{Arrieta2020,Yedjour2024}.
\\\\
Stability refers to the consistency of an explanation when subjected to different subsets of 
training data or the inherent stochasticity of the training process \cite{Ye2024}. Although 
stability has been studied primarily in image classification settings, \cite{Ye2024} establish that 
the underlying mechanism, explanation variance induced by model sensitivity to small perturbations, is a 
property of the model's loss landscape geometry rather than the data modality, and therefore generalises to 
tabular settings. In cancer genomics specifically, where batch effects and sample preparation variation 
introduce systematic noise into the feature space \cite{Partin2023}, an explanation that changes drastically 
due to negligible perturbations is considered fragile and clinically unreliable. This fragility not only erodes user 
trust \cite{Wicker2022} but also hinders the clinician's ability to discover genuine, generalisable biological 
relationships, restricting explainability to debugging misclassifications rather than generating medical insight \cite{Ye2024}. 

\subsubsection{Explanation Fragility and the Generalisation Gap}
As established in Section 3.4.2, the Clever Hans effect represents a core threat to explanation reliability in 
overfitted models. The critical question for this project is how to detect and quantify this fragility before 
explanations are extracted. The Generalisation Gap, the difference between training and test performance 
\cite{Keskar2016,Springer2024}, serves as this primary diagnostic indicator, as models with a 
wide gap tend to overfit to training artefacts rather than learning genuine biological structure. The danger of 
this overfitting is hidden by traditional interpretability; \cite{Adebayo2018}, through sanity checks, showed 
that popular XAI methods are independent of the model’s parameters. This means that they can produce explanations 
that visually look convincing even for models that have actually just memorised noise. Consequently, this highlights 
the necessity of shifting toward a 'RED XAI' approach \cite{Biecek2024}, using XAI as an active diagnostic tool to 
interrogate the model and ensure it is learning grounded biological truths rather than mathematical coincidences resulting 
from reality drift \cite{Longo2024}.

\subsubsection{Multi-Criteria Decision Analysis: PROMETHEE II}
Given that fidelity, comprehensibility, and stability often exist in a state of trade-off (e.g., a highly 
faithful rule may be too complex to understand), a structured method for ranking rule quality is required 
\cite{Yedjour2024,Ye2024,Brans1985}.
\\\\
PROMETHEE II (Preference Ranking Organisation Method for Enrichment Evaluations) is a Multi-Criteria Decision 
Analysis (MCDA) technique designed to rank alternatives based on multiple conflicting criteria 
\cite{Yedjour2024,Matejova2025}. Unlike simple averaging, PROMETHEE II uses pairwise comparisons 
to establish an outranking relationship between different explanations, producing a complete ranking of all 
alternatives rather than only identifying a subset of optimal solutions as methods such as TOPSIS or VIKOR do 
\cite{Yedjour2024,Matejova2025}. This complete ranking property is particularly important in 
the rule extraction context, where the goal is not merely to identify a single best rule set but to compare the 
full distribution of rule quality across the baseline and SAM-regularised models. Recent applications in rule 
extraction through the GPRO\_RX system have demonstrated that PROMETHEE-assisted ranking can successfully 
identify optimal rules that maintain high data coverage while significantly reducing complexity 
\cite{Yedjour2024}. Following this precedent, and given that GPRO\_RX represents the closest methodological 
antecedent to the evolutionary rule extraction pipeline used in this project, PROMETHEE II is adopted as the 
ranking framework here.  

\subsubsection{Biological Grounding: Pathway Enrichment (GSEA)}
Finally, in the domain of cancer genomics, statistical fidelity is insufficient without biological verification. 
Gene Set Enrichment Analysis (GSEA) is a method used to determine whether a set of genes (such as those 
identified in an extracted rule) shows statistically significant, concordant differences between two biological 
states (e.g., cancer vs. healthy) \cite{Subramanian2005}. By cross-referencing extracted rules with 
established databases like KEGG or Enrichr, researchers can verify if the AI’s "logic" aligns with known 
molecular pathways (e.g., Apoptosis or DNA repair), thereby providing the ultimate "sanity check" for clinical 
relevance.
